\documentclass{beamer}

% --- Configurazione della Lingua e dei Font ---
\usepackage[utf8]{inputenc}
\usepackage[english]{babel}
\usepackage{xcolor}
\usepackage[edges]{forest}

% --- Scelta del Tema (Personalizzabile) ---
\usetheme{Berkeley}
\usecolortheme{dolphin}

% --- Metadati della Presentazione ---
\bfseries\title{LLM Benchmark}
\author[A. Nome]{Nome del Relatore}
\date{\today}

% --- Inizio del Documento ---
\begin{document}

% 1. Slide del Titolo (Title page)
\begin{frame}
    \titlepage
\end{frame}

% 2. Slide dell'Indice (Sommario)
\begin{frame}{Contents}
    \tableofcontents
\end{frame}

\section{Introduction}
% 3. Slide Standard con Elenco Puntato ed Effetto Pausa
\begin{frame}[c]
    \centering
    
    % Testo molto grande (Huge), in grassetto (bfseries) e di colore blu elegante
    {\Huge\textcolor{blue}{Introduction}}
    
    \vspace{0.1cm} % Spazio elegante tra la parola e la linea di sottolineatura
    
    % Una linea sottile di sottolineatura (lunga 5cm e spessa 1.5pt)
    \textcolor{black}{\rule{7cm}{1pt}}
    
\end{frame}


\begin{frame}{Goal and Keypoints}
    \textcolor{blue}{Goal:} create a reliable, scalable and user-friendly benchmark to test the planning capabilities of LLMs on planning problems.
    \vspace{1cm}

    \textcolor{blue}{\textbf{Keypoints}}
    \begin{itemize}
        \item Problems Classification
        \item Hardware Platform and LLM Models Selected
        \item Outputs and Evaluation Metrics
    \end{itemize}
\end{frame}

\begin{frame}{General Structure}
        \begin{columns}
        % Prima Colonna
        \begin{column}{0.5\textwidth}\footnotesize
            This slide is used to show the structure of the repo, under the root folder we have:
            \begin{itemize}
                \item \textbf{analysis}: contains the material needed to write reports and the report itself.
                \item \textbf{archive}: contains material provided by the professor and a PDDL guide written to help understand key features of the PDDL.
                \item \textbf{Benchmark Framework}: main folder containing the work done to create this benchmark.
            \end{itemize}
        \end{column}
        
        % Seconda Colonna
        \begin{column}{0.5\textwidth}
            \begin{forest}
            for tree={
                font=\ttfamily,
                grow'=0,
                child anchor=west,
                parent anchor=south,
                anchor=west,
                calign=fixed edge angles,
                edge={draw, semithick},
                l sep = 7pt,       % Ridotto lo spazio orizzontale tra i livelli
                s sep = 3pt,       % Ridotto lo spazio verticale tra le cartelle
                inner sep = 2.5pt,  % ridotto spazio attorno al testo
                folder, % stile cartella
            }
            [LLM Benchmark/
                [analysis/
                [domain complexity/]
                [notebooks/]
                [Reports/]
                ]
                [archive/
                [Papers/]
                [Planning Domains/]
                [References/]
                [Quick PDDL Guide/]
                ]
                [Benchmark Framework/
                [...]
                ]
                [README.md]
            ]
            \end{forest}
        \end{column}
    \end{columns}
    
\end{frame}

\begin{frame}{General Structure II}
        \begin{columns}
        % Prima Colonna
        \begin{column}{0.5\textwidth}\scriptsize
            Now the Benchmark Framework folder is highlighted, note that to simplify not all the sub folders are shown.
            \begin{itemize}
                \setlength{\itemsep}{2pt}  % Spazio tra un punto e l'altro
                \setlength{\parskip}{0pt}  % Spazio tra paragrafi interni
                \setlength{\parsep}{0pt}   % Spazio di parsing
                \setlength{\topsep}{0pt}   % Spazio sopra l'elenco
                \setlength{\partopsep}{0pt}% Spazio extra sopra l'elenco
                
                \item \textit{evaluators}: code for running evaluation and parsing of oututs
                \item \textit{Leonardo script}: files to gain access to leonardo booster
                \item \textit{outputs}: store outputs of benchmark
                \item \textit{prompts}: prompts written to be feed to LLMs 
                \item \textit{runner}: contains code to run some tests with LLMs on the booster
                \item \textit{scripts}: contains code necessary for the scoring of the inputs
                \item \textit{tasks}: problems selected
                \item \textit{utils}: executable of validator
                \item \textit{results}: heatmaps showing output metric results
                \item \textit{run benchmark}: file that runs whole working sequence
            \end{itemize}
        \end{column}
        
        % Seconda Colonna
        \begin{column}{0.5\textwidth}
            \begin{forest}
            for tree={
                font=\ttfamily,
                grow'=0,
                child anchor=west,
                parent anchor=south,
                anchor=west,
                calign=fixed edge angles,
                edge={draw, semithick},
                l sep = 7pt,       % Ridotto lo spazio orizzontale tra i livelli
                s sep = 4pt,       % Ridotto lo spazio verticale tra le cartelle
                inner sep = 2.5pt,  % ridotto spazio attorno al testo
                folder, % stile cartella
            }
            [Benchmark Framework/
                [evaluators/]
                [Leonardo script/]
                [outputs/]
                [prompts/]
                [runner/]
                [scripts/]
                [tasks/]
                [utils/]
                [results/]
                [run benchmark.py]
            ]
            \end{forest}
        \end{column}
    \end{columns}
    
\end{frame}


\section{Problems Calssification}
\begin{frame}[c]
    \centering
    
    % Testo molto grande (Huge), in grassetto (bfseries) e di colore blu elegante
    {\Huge\textcolor{blue}{Problems' Classification}}
    
    \vspace{0.1cm} % Spazio elegante tra la parola e la linea di sottolineatura
    
    % Una linea sottile di sottolineatura (lunga 5cm e spessa 1.5pt)
    \textcolor{black}{\rule{10cm}{1pt}}
    
\end{frame}

\begin{frame}{Classification Purpose}
    \textcolor{blue}{\textbf{Purpose:}} In order to obtain reliable outputs to be evaluated after we for sure need a way to classify the PDDL problems to be fed to the LLMs chosen, since we had available 20 different problems a easy but effective way to classify them was needed to choose wisely the problems to be used.\\
    The main idea behind this Input Classification develops around two main concepts:
    \begin{itemize}
        \item  need of a numerical score representy the difficulty of the problem.
        \item classify problems based on the type of reasoning to be employed to solve the problem.
    \end{itemize}

\end{frame}

\begin{frame}{Implementation: Numerical Score}\footnotesize
    As stated before the idea of using a numerical score to classify the difficulty of the problems came up because a relatively simple and repeatable way to evaluate problems was needed; in order to obtain the score the code in \textit{LLM benchmark/scripts/score domains complexity.py} has been designed.\\
    The main idea behind the computation of the score is to evaluate for each instance of a PDDL problem the presence or the absence of features that we believed could increment or decrement the difficuly of a problem, specifically:
    \begin{itemize}
        \item N. of Objects
        \item N of possible Actions
        \item N. of Goals
        \item 
    \end{itemize}
    A raw score for each instance of a problem is therefore computed then for each domain a summary is obtained.

\end{frame}


\begin{frame}{Implementation: Numerical Score II}\footnotesize
    \begin{itemize}
        \item 
    \end{itemize}

\end{frame}


\begin{frame}{Implementation: Reasoning Type}
    
\end{frame}


%NEW SECTION
\section{Hardware platform and Models selected}

\begin{frame}[c]
    \centering
    
    % Testo molto grande (Huge), in grassetto (bfseries) e di colore blu elegante
    {\Huge\textcolor{blue}{Hardware platform and Models selected}}
    
    \vspace{0.1cm} % Spazio elegante tra la parola e la linea di sottolineatura
    
    % Una linea sottile di sottolineatura (lunga 5cm e spessa 1.5pt)
    \textcolor{black}{\rule{10cm}{1pt}}
    
\end{frame}

\begin{frame}{Hardware Platform and Models Selected}

    \Huge\textcolor{red}{DA RIEMPIRE}

\end{frame}

\begin{frame}{Hardware Platform and Models Selected II}

    \Huge\textcolor{red}{DA RIEMPIRE}

\end{frame}


% NEW SECTION
\section{Outputs and Evaluation Metrics}
\begin{frame}[c]
    \centering
    
    % Testo molto grande (Huge), in grassetto (bfseries) e di colore blu elegante
    {\Huge\textcolor{blue}{Output and Evaluation Metrics}}
    
    \vspace{0.1cm} % Spazio elegante tra la parola e la linea di sottolineatura
    
    % Una linea sottile di sottolineatura (lunga 5cm e spessa 1.5pt)
    \textcolor{black}{\rule{10cm}{1pt}}
    
\end{frame}


\begin{frame}{Output Structure}
    \Huge\textcolor{red}{DA RIEMPIRE}
\end{frame}
\begin{frame}{Evaluation Metrics}
    \Huge\textcolor{red}{DA RIEMPIRE}
\end{frame}


\section{Conclusion}
\begin{frame}[c]
    \centering
    
    % Testo molto grande (Huge), in grassetto (bfseries) e di colore blu elegante
    {\Huge\textcolor{blue}{Conclusion}}
    
    \vspace{0.1cm} % Spazio elegante tra la parola e la linea di sottolineatura
    
    % Una linea sottile di sottolineatura (lunga 5cm e spessa 1.5pt)
    \textcolor{black}{\rule{10cm}{1pt}}
    
\end{frame}

\begin{frame}{Conclusion}
    \Huge\textcolor{red}{DA RIEMPIRE}
\end{frame}



\end{document}